\section{Sensitivity to the Modeling Assumptions}
\label{sec:robustness}
% SWEVO expansion: tab:robust-final is back in the main text (it was moved to the supplement by editor cut C2 in the
% MPCE version); the full re-evaluation stays in the supplement (Table S-robust). tab:F-bench and tab:F-equiv are
% captured in 08_beyond.tex (\CaptureSuppTables{analysis/mpce_supp_direction.tex}).
The benchmark fixes a power curve, a wake model, a minimum spacing, a boundary rule and a discretization of the wind rose. We re-evaluated every feasible final layout of the main comparison, without re-optimization, under a cubic power curve, a Gaussian wake and finer direction bins (Table~\ref{tab:robust-final}; all methods: Table~\ref{S-tab:robust}); this asks whether the delivered layouts keep their order under another model, not which layouts the methods would find for it; only for the minimum spacing were the methods also re-optimized on benchmark cases (Section~\ref{sec:spacing}).
% >>>>> begin analysis/mpce_tab_robust_final.tex
%% generated by mpce_results.py -- do not edit by hand
\begin{table}[!t]
\centering
\caption{Robustness to the benchmark model: final layouts re-evaluated, not re-optimized.}
\label{tab:robust-final}
\scriptsize\setlength{\tabcolsep}{2.4pt}
\begin{tabular}{lcccccc}
\toprule
& \multicolumn{2}{c}{$\Delta$ (\%)} & & & & \\
\cmidrule(lr){2-3}
Model & DS I & DS II & Rank & Best & $\bar\tau$ & Same \\
\midrule
Benchmark & +0.0 & +0.0 & 1.82 & PSO-VNS & -- & -- \\
Cubic curve & $-17.7$ & $-36.9$ & 1.82 & PSO-VNS & 0.95 & 94 \\
Cubic, 25~m/s cut-out & $-21.5$ & $-37.0$ & 1.80 & PSO-VNS & 0.95 & 93 \\
Gaussian wake & $-0.3$ & $-0.4$ & 1.97 & PSO & 0.66 & 54 \\
\bottomrule
\end{tabular}
\par\vspace{2pt}\parbox{\columnwidth}{\scriptsize All feasible final layouts of the 68 cases. $\Delta$: mean relative change of the objective (\%); Rank: average rank of PSO-VNS; Best: best-ranked method; $\bar\tau$: mean Kendall correlation between the benchmark and the alternative ordering within a case; Same: cases (\%) with unchanged best method. Ranks of all methods: supplementary material.}
\end{table}
% <<<<< end analysis/mpce_tab_robust_final.tex

\subsection{Power Curve}
\label{sec:powercurve}
A cubic power curve between cut-in and rated speed~\cite{Carrillo2013} (also with a 25~m/s cut-out) lowers the expected power by 17.7\% (Data Set~I) and 36.9\% (Data Set~II), so our energy values are benchmark values, but it leaves the ranking intact ($\bar\tau=0.95$; same best method in 94\% of the cases; PSO-VNS first, also with the cut-out). This check concerns the rankings: the $\pm0.05$~pp equivalence verdicts (Sections~\ref{sec:psovns-pso} and~\ref{sec:ablation}) were evaluated under the linear benchmark curve only.
% CHECK-FINAL [C21]: robustness.Cubic/CubicCutout: tau >= 0.9, PSO-VNS rank position 1

\subsection{Wake Model}
\label{sec:gaussian}
The uncalibrated Gaussian wake of~\cite{Bastankhah2014} ($K_{\rm G}=0.04$; Section~\ref{S-sec:S-gauss}) differs from the Jensen cone ($K=0.075$) in its smooth lateral profile and expansion rate and, unlike the cone test of the benchmark, gives no deficit to turbines upstream of the wake source. The mean objective changes little (Table~\ref{tab:robust-final}), but the wake loss of the same final layouts differs between the two models by 0.38~pp on average (median 0.27~pp), 7.6 times the equivalence margin (Section~\ref{sec:setup-stats}), and by more than the margin for 86\% of the 15{,}614 layouts, whereas the per-case difference PSO-VNS${}-{}$PSO changes by only 0.031~pp on average. The ordering is less stable ($\bar\tau=0.66$), and PSO is narrowly first (1.91 vs.\ 1.97); with 1$^\circ$ direction bins, PSO-VNS is first under the Gaussian wake as well (Section~\ref{sec:direction}).
% CHECK-EXTRA [X30]: Jensen vs Gaussian wake loss of the same layouts: mean \NXModelShiftMean pp (median \NXModelShiftMedian; \NXModelShiftRatio x margin; > margin for \NXModelShiftAboveMarginPct % of \NXModelShiftNLayouts layouts); per-case PSO-VNS - PSO changes by \NXModelShiftPairMean pp on average (below the margin)
% CHECK-FINAL [C22]: robustness.Gauss (15 deg): tau < 0.8, PSOC first, PSOBV second (1.91 vs 1.97); CHECK-DIR [F09]: PSO-VNS first under the Gaussian wake at 1 deg

\subsection{Minimum Spacing}
\label{sec:spacing}
The optimized layouts often lie on the $4D$ constraint: 49\% of the feasible final layouts also satisfy $5D$ and 36\% $6D$ (4.1\% and 0.1\% for $N=11$--15, where the fixed site leaves less room), so they do not transfer to larger spacings without re-optimization (site capacity: Table~\ref{S-tab:capacity}). After the primary analysis, the eight methods and RSD-VNS were re-optimized at $5D$ and $6D$ on the 12 split cases (30 seeds, 6,030 evaluations from random starts; Section~\ref{S-sec:S-rev-spacing}); the two cases $r=500$~m, $N=10$ admit no layout at these spacings, so the same 10 cases are compared at every spacing. PSO-VNS keeps the best average rank (1.40, 1.60 and 1.90 at $4D$, $5D$ and $6D$), and its mean advantage over PSO is larger at $5D$ and $6D$ ($-0.094$ and $-0.068$~pp, the latter over the 8 cases in which PSO has at least 15 feasible runs) than at $4D$ ($-0.035$~pp), with 90\% intervals below zero at the seed and the case level at $5D$ and $6D$ (e.g., $5D$: [$-$0.135, $-$0.052] and [$-$0.200, $-$0.006]; Table~\ref{S-tab:S-rev-spacing}). MS-SLSQP moves up with the spacing (average rank 5.50, 4.30 and 3.00; second at $6D$). The SSA phase gives no gain over disc sampling at $5D$ (SSA-VNS${}-{}$RSD-VNS: $+0.043$~pp, case-level 90\% CI [$-$0.011, 0.103]), but at $6D$ SSA-VNS has the lower mean wake loss in all 8 cases in which both qualify ($-0.136$~pp, [$-$0.215, $-$0.073], wholly below $-0.05$~pp; Wilcoxon $p=0.0078$). The finding that salp-swarm first phases give no gain over random sampling in the farm (Section~\ref{sec:ablation}) thus holds at $4D$ and $5D$ but not at $6D$; with 8 to 10 cases, these results are estimates for the split cases, not general effects. An all-run paired outcome (Table~\ref{S-tab:S-allrun-spacing}; feasible beats infeasible) adds reliability to this picture: PSO-VNS is feasible in every run at $5D$ and in 90.0\% at $6D$, and its score against PSO is 0.593 and 0.613 (95\% case-bootstrap intervals [0.497, 0.693] and [0.548, 0.670]). At $6D$, MS-SLSQP is feasible more often (95.0\%), and the all-run score of PSO-VNS against it, 0.580 [0.432, 0.733], does not show higher reliability, although PSO-VNS has the better conditional-quality rank.
% CHECK-FINAL [C23]: spacing_share.n11_15.pct_5d < 10 ("often lie on the 4D constraint")
% verified (manual): shares 49 % / 36 % (all N) vs 4.1 % / 0.1 % (N = 11-15) from mpce_numbers.tex (\NSp...); constructive capacity at 4D/5D/6D: Table S-capacity

\subsection{Boundary Rule}
\label{sec:boundary}
All methods clip the coordinates to the bounding square and leave the circle to the penalty (Section~\ref{sec:constraints}). The rule matters: an SSA that projects turbines radially onto the circle instead attains higher mean objectives than the original SSA in all six representative cases ($p\le0.0036$), as the projection returns a turbine to the site directly. All methods therefore use the same rule: the comparison is one of search strategies under a common repair, which is not thereby neutral across methods (Section~\ref{sec:discussion}).
% CHECK-FINAL [C41]: boundary_rule n_higher == n_cases == 6 and max_p < 0.01

\subsection{Direction Resolution}
\label{sec:direction}
The benchmark evaluates the wake only at the centers of 24 direction bins of 15$^\circ$, so a pair of turbines can be placed with no bin center in its narrow wake interval (Section~\ref{sec:model}). We split each bin into 3 (5$^\circ$) or 15 (1$^\circ$) sub-bins with its Weibull parameters and an equal share of its frequency (wake-free power unchanged) and re-evaluated all 21{,}716 feasible final layouts (Table~\ref{S-tab:F-bench}).
% CHECK-DIR [F01]: one sub-bin reproduces the benchmark objective; wake-free objective identical for 1, 3 and 15 sub-bins
The 1$^\circ$ rose raises the mean wake loss of every method, by 0.40~pp (DE) to 1.04~pp (PSO; PSO-VNS: 1.322\% to 2.335\%), mostly already with 5$^\circ$ sub-bins: optimized layouts exploit the gaps between the coarse bins, those of PSO and the VNS-based methods most. PSO-VNS still ranks first (2.09; PSO 2.37, the only method not significantly different), also under the Gaussian wake at 1$^\circ$ (2.22 vs.\ 2.35); PSO's narrow lead under the Gaussian wake with 15$^\circ$ bins (1.91 vs.\ 1.97; $\bar\tau=0.66$) is a property of the coarse rose. Beyond the top two the order changes ($\bar\tau=0.52$; same best method in 37\% of the cases): MS-SLSQP moves from sixth to third, and its gap to PSO-VNS shrinks from 0.435 to 0.135~pp but remains significant.
% CHECK-DIR [F03]: loss rise of every method, PSO-VNS 1.322 -> 2.335; [F04]: most of the change already with 5-deg sub-bins (> 70 % for PSO-VNS); [F05]: VNS-based methods and PSO rise most, DE/MS-SLSQP least ("VNS-based methods most"); [F06]: PSO-VNS first at 5 and 1 deg, PSO the only one not significantly different (Holm); [F09]: Gaussian: PSO first at 15 deg, PSO-VNS first at 1 deg; [F07]: MS-SLSQP 6th -> 3rd, tau 0.52, same best 37 %; [F08]: MS-SLSQP gap shrinks from 0.435 to 0.135 pp, still significant (Holm)
% Page budget: kept in the supplementary material: \SuppTable{tab:F-bench}

For PSO-VNS vs.\ PSO (Table~\ref{S-tab:F-bench}; details: Table~\ref{S-tab:F-equiv}), the case-mean difference becomes \ensuremath {-}0.043~pp (90\% CI [\ensuremath {-}0.062, \ensuremath {-}0.025]): a small but significant advantage ($p=\ensuremath {2.5\times 10^{-4}}$), not equivalence at $\pm0.05$~pp (minimal margin 0.062~pp); the same holds with 5$^\circ$ sub-bins. The sign of the per-case difference changes in 26 of the 60 cases with a nonzero difference, so the per-case winner often depends on the discretization. On the 44 non-trivial cases, equivalence is shown at neither resolution, and the Data Set~II $N\ge10$ advantage remains (10 of 10 cases). The component-analysis conclusions hold, but the small SSA gain over RS-VNS is no longer significant (Table~\ref{S-tab:F-abl}). The order beyond the top two and the equivalence verdict are thus properties of the discretized benchmark; the first place of PSO-VNS, as a ranking of re-evaluated layouts, is not. Which method would lead if the search itself used 1$^\circ$ bins was not tested.
% CHECK-DIR [F10]: PSO-VNS - PSO -0.043, CI [-0.062, -0.025], not equivalent, p < 0.05; [F11]: same with 5-deg sub-bins; [F12]: per-case sign changes in 26 of 60 cases; [F13]: non-trivial cases: equivalence at neither resolution; [F14]: DS II N >= 10 10/10; [F16]: component-analysis conclusions hold (SSA-VNS and LX-SSA-VNS worse than RSD-VNS, PSO-VNS better than RSD-VNS, RS-VNS and VNS, PSO-VNS best of nine); [F17]: SSA-VNS vs RS-VNS not significant at 1 deg
% Over the page budget (+0.5 pages in the private build), kept in the supplement: \SuppTable{tab:F-equiv}
